\section{Exercice de conception : Plan de salle}
On souhaite écrire un programme qui aide un enseignant à représenter le plan d'une salle de classe, organisée en rangées de tables.

\subsection{Question 1 : Liste des rangées}

\begin{UPSTIManipulation}{Plan de salle -- Liste des rangées}
	Écrire un programme qui :
	\begin{enumerate}
		\item Demande à l'utilisateur le nombre de rangées \texttt{n} de la salle (\texttt{get\_int}).
		\item Affiche, pour chaque rangée numérotée de 1 à \texttt{n}, une ligne de la forme \texttt{Rangee <numero> : libre}.
	\end{enumerate}

	Exemple d'exécution :
	\begin{lstlisting}[language=bash,style=console]
Entrez le nombre de rangees :
3
Rangee 1 : libre
Rangee 2 : libre
Rangee 3 : libre
\end{lstlisting}
\end{UPSTIManipulation}

\subsection{Question 2 : Plan complet de la salle}

\begin{UPSTIManipulation}{Plan de salle -- Plan complet}
	On souhaite maintenant afficher un plan visuel de la salle entière.

	Écrire un programme qui :
	\begin{enumerate}
		\item Demande à l'utilisateur le nombre de rangées \texttt{n} de la salle (\texttt{get\_int}).
		\item Demande à l'utilisateur le nombre de places \texttt{m} par rangée (\texttt{get\_int}).
		\item Affiche le plan de la salle : \texttt{n} lignes de \texttt{m} places chacune, chaque place étant représentée par le symbole \texttt{o}, séparé du suivant par un espace.
	\end{enumerate}

	Exemple d'exécution :
	\begin{lstlisting}[language=bash,style=console]
Entrez le nombre de rangees :
3
Entrez le nombre de places par rangee :
4
o o o o
o o o o
o o o o
\end{lstlisting}
\end{UPSTIManipulation}

\subsection{Pour aller plus loin (optionnel)}

\begin{UPSTIManipulation}{Plan de salle -- Rangées occupées}
	On souhaite maintenant que les rangées de numéro \textbf{impair} soient occupées.

	Modifier les deux programmes précédents :
	\begin{enumerate}
		\item Dans le programme de la Question 1 : afficher \texttt{Rangee <numero> : occupee} pour une rangée de numéro impair, et \texttt{Rangee <numero> : libre} pour une rangée de numéro pair.
		\item Dans le programme de la Question 2 : afficher le symbole \texttt{x} pour toutes les places d'une rangée de numéro impair, et le symbole \texttt{o} pour toutes les places d'une rangée de numéro pair.
	\end{enumerate}

	Exemple d'exécution (programme de la Question 2, avec \texttt{n = 3} et \texttt{m = 4}) :
	\begin{lstlisting}[language=bash,style=console]
Entrez le nombre de rangees :
3
Entrez le nombre de places par rangee :
4
x x x x
o o o o
x x x x
\end{lstlisting}
\end{UPSTIManipulation}


\begin{UPSTIManipulation}
		      \begin{lstlisting}[language=bash,style=console]
submit50 iut-geii-annecy/exercices/2026/info1/tp_autoeval_boucles/plan_salle/liste_rangees
submit50 iut-geii-annecy/exercices/2026/info1/tp_autoeval_boucles/plan_salle/plan_complet
\end{lstlisting}
	
\end{UPSTIManipulation}